\section{Fear Brought Before the Savior}

Fear of hell can interrupt the movement toward sin. When a desired pleasure fills the imagination and love has become faint, remembrance of the end can give the will a reason to refuse. St.~Ignatius of Loyola expressly seeks this help in the First Week of the \emph{Spiritual Exercises}. The request introducing his meditation on hell asks for an interior sense of its pains so that, if faults cause the exercitant to forget the love of the eternal Lord, fear may help keep him from sin (nos.~65--71).

Ignatius asks for fear's assistance when the love of the Lord has been forgotten through one's faults. The exercitant knows the evil of sin, but its attraction may occupy his imagination more powerfully than its consequences. He needs to recognize inwardly what his judgment already acknowledges. Remembering punishment can interrupt his willing pursuit of a pleasure that has seemed harmless because its end remained distant. The request seeks help for that moment of weakness, when a truth confessed in general must become a reason for refusing a particular act.

His method engages the imagination through the five senses. The exercitant considers the fire, the cries, the bitterness, and the other pains presented in the exercise. Sight is not left to wander after the object of temptation; hearing, taste, smell, and touch are drawn into the consideration of its end. The method presses against the habit of imagining sin's attraction vividly while keeping its consequences abstract.

These are images used in meditation, not a new revelation of the physical arrangement of hell. Ignatius asks the exercitant to consider the horror of condemnation before the Savior. A person who can profitably use the senses in this way should keep the exercise directed toward its intended fruit. Another may be held more deeply by Christ's words of departure, by the permanence of the loss, or by remembrance of mercy. In each case the will is to become more ready to refuse sin and belong to God.

Ignatius ends with a colloquy addressed to Christ. The exercitant recalls those who did not believe and those who believed but did not obey, then gives thanks that his own life has been preserved. He remembers how the Lord has continued to show him pity and mercy. The remembrance of punishment thus brings him to a personal acknowledgment of what he has received. He has been given time to return to Christ, despite faults for which he can claim no excuse. Gratitude makes his continued refusal still harder to defend.

The colloquy also gives the frightened sinner someone to whom he can turn. Recognizing the danger does not itself supply the strength to resist temptation. In prayer he brings his weakness and past resistance to the Lord who bore the Passion for sinners. A resolution dependent only upon his present feeling will weaken when the feeling passes. Recourse to Christ asks for help during the temptation itself, when the attraction returns and the warning again becomes easy to dismiss. Ignatius's meditation ends in that dependence upon the Savior.

Fear also needs truth about its own object. The evil to be feared is sin and final separation from God. God himself is good, and his goodness makes that separation dreadful. A fear that drives a person to conceal sin or flee the sacrament has lost its proper direction.

As charity grows, sorrow becomes more deeply concerned with God himself. The soul grieves that it has offended such goodness and refused such love. Fear need not be disowned for this to happen. It can continue to guard fidelity while love gives a fuller motive for it. Polycarp's confession shows the union: he knows the eternal danger, and he refuses to betray his King and Savior.

The fruit of this meditation can be considered in the act which follows it. A quiet refusal of a familiar temptation may require more fidelity than a strong emotion during prayer. The will has to obey when the object again appears attractive and the imagined consequences seem remote. Gratitude to Christ can help sustain that obedience, because the penitent remembers whose mercy he is receiving. Intense feeling may come or remain absent. What needs attention afterward is the willingness to use the help offered when the occasion for sin returns.

Ignatius's surrounding instruction on general confession also preserves proportion. He proposes it to those who choose it freely, explains its possible benefits, and states that one who confesses annually is not thereby obliged to make a general confession (no.~44). An annual preparation for death should therefore be undertaken with sound counsel. Its purpose is a sincere return and a more faithful life.

